\section{Discussion}\label{sec:discussion}

\subsection{Design principles}

Design science research should leave behind knowledge that transfers beyond the instantiated artifact \citep{hevner2004}. We condense the evaluation into six design principles (DPs) for affective components of streaming video information systems. Each principle is tied to the evidence that supports it.

\begin{description}
\item[DP1.] \emph{Fuse late, and choose the first stage by measured cost.} Late fusion of separately trained audio and visual foundation features gave the largest accuracy gain in the study (+10.5 points). It also kept the modalities separable for gating and explanation. Which model runs first should follow from measured cost and accuracy, not from model size. Once face detection was counted, the speech transformer was cheaper than the ``cheap'' visual pipeline. A selection made inside the training folds chose it every time (Table~\ref{tab:orders}).
\item[DP2.] \emph{Train for the degradations you can anticipate.} Multi-condition training was the most effective single measure for robustness. It recovered most of the accuracy lost to noise and compression, extended to unseen noise levels, and cost nothing at inference (Table~\ref{tab:mct}). It did not cover every shift: low light was not helped.
\item[DP3.] \emph{Do not use output confidence as a health signal.} With clean-trained models, confidence rose while accuracy collapsed under audio noise. Confidence-gated escalation and confidence-based drift detection then failed together (Tables~\ref{tab:robust} and \ref{tab:drift}). After robust training, confidence no longer rose, but calibration still deteriorated under shift. Confidence remains useful for allocating compute on in-distribution data. It needs an independent check that the inputs are still in distribution.
\item[DP4.] \emph{Monitor every channel's inputs, at the stream level, against two references.} A distance-to-training score on cheap per-channel features, fed to an adaptive-window detector, detected every degradation and attributed it to the right channel. Measured against clean data, it tells operators that and where the input has changed. Measured against the training envelope, it tells the system whether the models can still be trusted on that channel. Per-clip thresholds were unreliable (Fig.~\ref{fig:monitors}). A monitor that observes only a fraction of a channel's clips detected degradations late and missed their end (Table~\ref{tab:integrated}). Natural speaker changes also triggered alarms, so reference statistics should be adapted to the speaker where possible.
\item[DP5.] \emph{Gate only on shifts outside the training envelope, and hand over to humans when no channel is trusted.} With clean-trained models, gating on stream-level alarms recovered most of the accuracy lost to audio noise. With robust models, gating on any departure from clean data was harmful, while gating on envelope alarms was safe. When both channels are degraded, no automatic fix is available, and the component should report the problem and mark its estimates as unreliable.
\item[DP6.] \emph{Recalibrate and audit fused outputs per group.} Product-rule fusion was over-confident, and the over-confidence differed by group. Recalibration fixed this for the full model but worsened it for one group in the cheap stage. Accuracy differences between groups were not fixed by calibration, and relative to the strongest single modality they were larger after fusion. Both have to be measured and reported as part of deployment.
\end{description}

\subsection{Implications for research}

For affective computing, our results suggest that benchmark accuracy is an incomplete criterion for components that will run on live streams. The strongest single model on clean data, wav2vec~2.0, was also the most fragile under noise. Confident errors under degradation do not show up in clean-data evaluation. Reporting accuracy under controlled channel degradations, together with the behaviour of confidence, would make results more informative for system builders. Zero-shot CLIP did not recognise acted facial emotion, although a classifier trained on its image features did. General-purpose vision--language models should therefore not be assumed to provide affect recognition without adaptation.

For IS research on analytics and AI management, the study offers an instance of what \citet{berente2021} call the management of autonomy and inscrutability. The artifact keeps a human in the loop where automatic correction runs out. It reports which channel is affected and since when. And it makes each decision decomposable into modality contributions. The design cycles also illustrate why IS artifacts that embed machine learning should be evaluated under the conditions of use, not only under the conditions of training \citep{shmueli2011}. Our first design worked on clean data and failed on degraded streams, and only the stream evaluation showed this. Our third design worked on degraded streams but became harmful once the models were trained robustly. The value of a component depends on the other components it is combined with.

\subsection{Implications for practice}

Organisations that run affect analytics on video, in customer service, tele-health, advertising research or live commerce, can take four practical points from this study. First, compute budgets are dominated by decoding, face detection and the largest models. Measuring the cost of each stage on the target hardware is cheap and changes design decisions. Second, training on realistic degraded copies of the data is the cheapest robustness measure, and a model's confidence should not be used as a quality-control signal for the stream it is analysing. Cheap per-channel input monitors do this job better, cost under a millisecond per clip once their features exist, and should observe every clip. Third, when a channel is flagged, an affect pipeline can keep running on the other channel with a predictable loss of accuracy. When both are flagged it should pause or label its output as unreliable. Fourth, group-level accuracy and calibration should be reported as part of model documentation. In the European Union, emotion recognition in workplaces and educational settings is prohibited except for medical or safety purposes, and other uses are high-risk and carry obligations for risk management, data governance, logging and human oversight \citep{euaiact2024}. The monitoring, logging and audit functions of the artifact map naturally onto such obligations. They do not make a prohibited use acceptable.

\subsection{Ethical considerations}

Automated inference of emotion raises concerns that technical safeguards do not remove. Facial and vocal expressions are not reliable read-outs of internal states \citep{barrett2019}. Our models estimate the emotion that actors were asked to \emph{portray}, and their outputs should be read as estimates of expressed affect. Affect analytics can be used to manipulate customers or to surveil employees. The design principles above concern the reliability of a component and do not settle whether a given use is legitimate. We used a public corpus whose actors consented to research use, and we release no new personal data.

\subsection{Limitations and future research}

The study has limitations that suggest further work. First, RAVDESS contains acted, frontal, studio-quality expressions by 24 North American actors speaking two neutral sentences. Absolute accuracies will be lower on spontaneous, in-the-wild data, and the language channel, which our design omits, carries much affect in natural conversation \citep{zadeh2018}. The relative findings, namely that confidence rises under noise and that input monitors detect channel degradation, should be tested on spontaneous corpora such as CMU-MOSEI and on cross-corpus shifts. Second, our degradations are synthetic and stationary within a segment, and the multi-condition training used two of the same synthetic degradation types that were tested. Its generalisation to real recording conditions is therefore likely to be overstated. Real streams degrade gradually, intermittently and in combination, and the detector settings (check interval, sensitivity, recovery asymmetry) will need tuning for such patterns. Third, the false-alarm results rest on replays of 1440 clips from 24 actors, and the speaker-blocked streams showed that monitors react to speaker changes. False-alarm rates on long real streams with many speakers remain to be measured. Fourth, we used light classifiers on frozen features. Fine-tuned or jointly trained models may be more accurate, but they may also behave differently under degradation. Fifth, the group analysis covers one binary attribute with 12 actors per group and cannot separate sex from individual expressive style. Larger and more diverse samples are needed before conclusions about demographic bias can be drawn. Sixth, cost was measured on one CPU type, and the relative costs of decoding, face detection and models will differ on GPUs and edge accelerators. The cost model, however, is easily recomputed for any hardware. Finally, we evaluated the artifact technically and did not study how operators use alarms and attributions. A field evaluation with practitioners, for example in a contact centre or tele-health setting, would complete the design science cycle.
